% ============================================================
%  SOAL + PEMBAHASAN  --  FISIKA  --  SSU-ITB 2026
%  SM-SSU ITB 2025  (nomor 1--30) | Paket 4
% ============================================================

\documentclass[11pt]{article}

\usepackage[utf8]{inputenc}
\usepackage[T1]{fontenc}
\usepackage[a4paper,margin=2.5cm,top=2.8cm]{geometry}
\usepackage{amsmath,amssymb}
\usepackage{graphicx}
\usepackage{enumitem}
\usepackage{multicol}
\usepackage{xcolor}
\usepackage[most]{tcolorbox}
\usepackage{fancyhdr}
\usepackage{booktabs}
\usepackage{icomma}
\usepackage{array}

% --- Warna Fisika ---
\definecolor{mapelcolor}{RGB}{20,70,150}
\definecolor{mapelbg}{RGB}{230,238,255}
\definecolor{pembcolor}{RGB}{100,50,130}
\definecolor{pembbg}{RGB}{245,238,252}
\definecolor{gold}{RGB}{210,160,30}
\definecolor{catatan}{RGB}{180,120,0}

\pagestyle{fancy}\fancyhf{}
\lhead{\small SSU-ITB 2026 --- Pembahasan}
\rhead{\small Paket 4}
\cfoot{\small\thepage}
\renewcommand{\headrulewidth}{0.4pt}

\newtcolorbox{soal}[1]{
  breakable,
  colback=mapelbg, colframe=mapelcolor,
  coltitle=white, fonttitle=\bfseries\sffamily,
  title=Nomor #1,
  boxrule=0.8pt, arc=3pt,
  left=3mm, right=3mm, top=2mm, bottom=2mm}

\newtcolorbox{pembox}{
  breakable,
  colback=pembbg, colframe=pembcolor,
  coltitle=white, fonttitle=\bfseries\sffamily,
  title=Pembahasan,
  boxrule=0.8pt, arc=3pt,
  left=3mm, right=3mm, top=2mm, bottom=2mm}

\newtcolorbox{catatbox}{
  breakable,
  colback=gold!10, colframe=catatan,
  coltitle=black, fonttitle=\bfseries\sffamily,
  title=Catatan,
  boxrule=0.8pt, arc=3pt,
  left=3mm, right=3mm, top=1.5mm, bottom=1.5mm}

\newcommand{\jawaban}[1]{%
  \noindent\colorbox{gold!30}{\parbox{\dimexpr\linewidth-2\fboxsep\relax}{%
    \textbf{Jawaban:} #1}}\par\smallskip}

\newcommand{\konsep}{\noindent\textit{Konsep:}\ \ignorespaces}

\newenvironment{pil}{%
  \begin{enumerate}[label=(\alph*),leftmargin=*,itemsep=1pt,topsep=3pt]}%
  {\end{enumerate}}
\newenvironment{pilB}{%
  \par\smallskip\begin{multicols}{2}
  \begin{enumerate}[label=(\alph*),leftmargin=*,itemsep=1pt,topsep=3pt]}%
  {\end{enumerate}\end{multicols}}

\newcommand{\gambar}[2][0.52\textwidth]{%
  \begin{center}\includegraphics[width=#1]{#2}\end{center}}

\linespread{1.05}

% ===================================================================
\begin{document}
\thispagestyle{empty}

\begin{center}
  \fbox{\parbox{0.6\textwidth}{\centering\bfseries\large
    LEMBAR SOAL DAN PEMBAHASAN}}
\end{center}
\vspace{0.5cm}

\begin{center}
  {\LARGE\bfseries FISIKA}
\end{center}
\vspace{0.5cm}

\begin{center}
  \begin{tabular}{@{\hspace{2cm}}ll@{\hspace{0.4cm}:@{\hspace{0.4cm}}}l}
    Jenjang      & & SMA / Sederajat               \\[0.3em]
    Hari/Tanggal & & \underline{\hspace{5cm}}      \\[0.3em]
    Waktu        & & \underline{\hspace{5cm}}      \\[0.3em]
    Paket Soal   & & 4                             \\
  \end{tabular}
\end{center}

\vspace{0.6cm}
\noindent\rule{\linewidth}{0.6pt}
\vspace{0.3cm}

\noindent\textbf{\underline{PETUNJUK UMUM}}
\begin{enumerate}[leftmargin=*,itemsep=4pt,topsep=4pt]
  \item Anda \textbf{tidak} diperbolehkan menggunakan sumber-sumber eksternal,
        termasuk internet, \textit{large language model} seperti ChatGPT, Claude,
        LLaMA, Gemini, Meta AI, dan sebagainya. Anda sangat dianjurkan untuk
        mencoba mengerjakan sendiri terlebih dahulu karena evaluasi ini dibuat
        untuk \textbf{mengukur} tingkat kepahaman; nilai $<70$ mengindikasikan
        \textbf{tidak lulus}, dan jika sudah melewati \textit{deadline} maka
        dianggap tidak mengerjakan yang berarti \textbf{tidak lulus}.

  \item Dokumen ini merupakan \textbf{lembar soal sekaligus pembahasan}; gunakan
        pada tahap \textbf{Review}. Di tahap ini Anda diharapkan mengerjakan ulang
        dengan disertakan penjelasan/pembelaan mengapa revisi jawaban adalah pilihan
        tersebut, dan tetap mencantumkan penjelasan pada soal yang walaupun pada
        penilaian sudah benar.

  \item Anda dianjurkan untuk mencantumkan caranya walaupun tidak termasuk
        penilaian, mengingat sifatnya adalah sebagai bahan pembelajaran.
        Mohon bertanggung jawab atas pekerjaan Anda sendiri.

  \item Terdapat tiga jenis soal: \textbf{Pilihan Ganda Biasa} (5 opsi, pilih 1);
        \textbf{Pilihan Ganda Majemuk Kompleks} (tabel \textbf{Benar}/\textbf{Salah},
        tiap pernyataan dinilai mandiri); dan \textbf{Isian Singkat}
        (tulis nilai/ekspresi).

  \item Soal berjumlah \textbf{30 butir}: nomor 1--30 dari latihan SM-SSU
        ITB 2025.

  \item Isilah detail informasi berikut.\\[0.3em]
        \textbf{Nama}\hspace{0.45cm}:\enspace\underline{\hspace{8cm}}
\end{enumerate}

\vspace{0.6cm}
\noindent\textit{Dengan menandatangani kolom di bawah ini, saya menyatakan telah membaca,
memahami, dan menyetujui seluruh peraturan yang tercantum di atas.}

\vspace{0.6cm}
\noindent\fbox{\parbox{5cm}{\vspace{2.2cm}\centering\small Tanda Tangan Peserta\vspace{0.3cm}}}

\clearpage

% ===================================================================
%  NOMOR 1--10
% ===================================================================

\begin{soal}{1}
\gambar[0.32\textwidth]{images/4/4-n1.png}
Sebuah roda silinder pejal berporos di $O$. Gaya $F_1$ dan $F_2$ bekerja
selama 20 s. Setelah $F_2$ dihilangkan, roda berhenti 20 s kemudian.
Perbandingan $F_2/F_1$ adalah \ldots
\begin{pilB}
  \item $3:2$ \item $2:1$ \item $1:3$ \item $2:3$ \item $1:2$
\end{pilB}
\end{soal}
\begin{pembox}
\jawaban{(b) $F_2/F_1 = 2:1$}
\konsep Torsi tangensial $\tau=FR$; percepatan sudut $\alpha=\tau_{\text{net}}/I$.

Selama 20 s pertama, $F_2$ dan $F_1$ berlawanan arah torsi:
$\alpha_1=\dfrac{(F_2-F_1)R}{I}$; kecepatan sudut akhir
$\omega=\alpha_1\times20=\dfrac{(F_2-F_1)R}{I}\cdot20$.

Setelah $F_2$ hilang, hanya $F_1$ yang melambatkan:
$\alpha_2=\dfrac{F_1 R}{I}$.
Roda berhenti setelah 20 s, sehingga $\omega=\alpha_2\times20$:
\[
  \frac{(F_2-F_1)R}{I}\cdot20=\frac{F_1 R}{I}\cdot20
  \;\Rightarrow\; F_2-F_1=F_1
  \;\Rightarrow\; F_2=2F_1.
\]
\end{pembox}

\begin{soal}{2}
Pegas digantung vertikal; beban $M$ gram bertambah panjang $\Delta l$.
Frekuensi pegas = 11× frekuensi bandul panjang $L$ dan massa $0{,}5M$.
Tentukan kebenaran setiap pernyataan!

\par\smallskip
{\renewcommand{\arraystretch}{1.4}%
\begin{tabular}{@{}p{0.76\linewidth}cc@{}}
\hline\textbf{Pernyataan}&\textbf{Benar}&\textbf{Salah}\\\hline
(1)~$f_s=\dfrac{1}{2\pi}\sqrt{\dfrac{g}{\Delta l}}$ (massa tidak berpengaruh).&$\bigcirc$&$\bigcirc$\\
(2)~$f_b=\dfrac{1}{2\pi}\sqrt{\dfrac{g}{L}}$ tidak bergantung massa.&$\bigcirc$&$\bigcirc$\\
(3)~Dari $f_s=11f_b$: $L=11\,\Delta l$.&$\bigcirc$&$\bigcirc$\\
(4)~Jika $\Delta l=1{,}9$ cm, maka $L=229{,}9$ cm.&$\bigcirc$&$\bigcirc$\\\hline
\end{tabular}}
\end{soal}
\begin{pembox}
\jawaban{(1)~\textbf{B}\quad(2)~\textbf{B}\quad(3)~\textbf{S}\quad(4)~\textbf{B}}
Untuk pegas vertikal: $k\Delta l=Mg\Rightarrow k/M=g/\Delta l$, sehingga
$f_s=\dfrac{1}{2\pi}\sqrt{k/M}=\dfrac{1}{2\pi}\sqrt{g/\Delta l}$
(tidak bergantung massa beban).
\begin{enumerate}[leftmargin=*,topsep=2pt,itemsep=2pt]
\item \textbf{Benar.} Formula $f_s$ benar.
\item \textbf{Benar.} Frekuensi bandul hanya bergantung $L$ dan $g$.
\item \textbf{Salah.} $f_s=11f_b\Rightarrow\sqrt{\tfrac{g}{\Delta l}}=11\sqrt{\tfrac{g}{L}}\Rightarrow L=121\,\Delta l$, \emph{bukan} $11\,\Delta l$.
\item \textbf{Benar.} $L=121\times1{,}9=229{,}9$ cm.
\end{enumerate}
\end{pembox}

\begin{soal}{3}
Gelombang sinusoidal ke arah negatif $x$: amplitudo 10 cm, $\lambda=314$ cm,
$f=2$ Hz, $y(x{=}10\text{ cm},t{=}0)=0$. Persamaannya \ldots
\begin{pil}
  \item $y=0{,}1\sin(2x-4\pi t-0{,}2)$
  \item $y=0{,}1\sin(2x+4\pi t)$
  \item $y=0{,}1\sin(2x+4\pi t-0{,}2)$
  \item $y=0{,}1\sin(2x+4\pi t+0{,}2)$
  \item $y=0{,}1\sin(2x-4\pi t)$
\end{pil}
\end{soal}
\begin{pembox}
\jawaban{(c) $y=0{,}1\sin(2x+4\pi t-0{,}2)$}
Gelombang ke $-x$: $y=A\sin(kx+\omega t+\phi)$.
$k=\dfrac{2\pi}{\lambda}=\dfrac{2\pi}{3{,}14}\approx2$ rad/m;\quad
$\omega=2\pi f=4\pi$ rad/s.

Syarat $y(0{,}1,0)=0$: $\sin(2\times0{,}1+\phi)=\sin(0{,}2+\phi)=0$,
sehingga $\phi=-0{,}2$.
\end{pembox}

\begin{soal}{4}
\gambar[0.3\textwidth]{images/4/4-2.png}
Balok $A$ (massa $M$, ketinggian $H$) menumbuk benda $B$ (massa $m$, diam)
di titik terendah, menempel. Ketinggian maksimum setelah tumbukan \ldots $H$.
\begin{pilB}
  \item $\!\left(\dfrac{M+m}{m}\right)^{\!2}$
  \item $\!\left(\dfrac{M}{M+m}\right)^{\!2}$
  \item $\!\left(\dfrac{m}{M+m}\right)^{\!2}$
  \item $\!\left(\dfrac{M}{M+m}\right)$
  \item $\!\left(\dfrac{M+m}{M}\right)^{\!2}$
\end{pilB}
\end{soal}
\begin{pembox}
\jawaban{(b) $\left(\dfrac{M}{M+m}\right)^{\!2}\!H$}
Sebelum tumbukan: $MgH=\tfrac12Mv^2\Rightarrow v=\sqrt{2gH}$.

Tumbukan menempel (momentum kekal):
$Mv=(M+m)v'\Rightarrow v'=\dfrac{M}{M+m}v$.

Setelah tumbukan: $h=\dfrac{v'^2}{2g}=\left(\dfrac{M}{M+m}\right)^2\cdot\dfrac{v^2}{2g}=\left(\dfrac{M}{M+m}\right)^2 H$.
\end{pembox}

\begin{soal}{5}
Apa pengaruh angin terhadap efek Doppler pada suara di udara?
\begin{pil}
  \item Efek Doppler terjadi selama ada angin meskipun sumber dan pengamat diam.
  \item Angin dapat mengubah besar efek Doppler, tetapi tidak dapat menyebabkannya.
  \item Angin dapat menyebabkan efek Doppler seperti sumber yang bergerak.
  \item Angin tidak memengaruhi efek Doppler sama sekali.
  \item Efek Doppler hanya terjadi jika angin bertiup ke arah pengamat.
\end{pil}
\end{soal}
\begin{pembox}
\jawaban{(b)}
Efek Doppler terjadi karena gerak relatif sumber/pengamat terhadap medium (udara).
Angin seragam yang menggerakkan seluruh medium tidak menciptakan gerak relatif
antara sumber dan medium, sehingga angin \emph{saja} tidak menimbulkan efek Doppler.
Namun jika sudah ada efek Doppler (sumber/pengamat bergerak), angin mengubah
kecepatan relatif terhadap medium sehingga besar efek Doppler berubah.
\end{pembox}

\begin{soal}{6}
Balok di bidang kasar ditarik gaya $F$ arah $0^\circ<\alpha<80^\circ$
terhadap perpindahan sehingga bergerak. Tentukan kebenaran setiap pernyataan!

\par\smallskip
{\renewcommand{\arraystretch}{1.4}%
\begin{tabular}{@{}p{0.76\linewidth}cc@{}}
\hline\textbf{Pernyataan}&\textbf{Benar}&\textbf{Salah}\\\hline
(1)~Kerja oleh gaya gesek pada balok bertanda positif.&$\bigcirc$&$\bigcirc$\\
(2)~Peningkatan energi kinetik balok sebanding dengan perpindahan.&$\bigcirc$&$\bigcirc$\\
(3)~Jika $\alpha$ berubah, hanya kerja dari $F$ yang terpengaruh.&$\bigcirc$&$\bigcirc$\\
(4)~Kerja gaya gravitasi dan gaya normal masing-masing bernilai nol.&$\bigcirc$&$\bigcirc$\\\hline
\end{tabular}}
\end{soal}
\begin{pembox}
\jawaban{(1)~\textbf{S}\quad(2)~\textbf{B}\quad(3)~\textbf{S}\quad(4)~\textbf{B}}
\begin{enumerate}[leftmargin=*,topsep=2pt,itemsep=2pt]
\item \textbf{Salah.} Gaya gesek berlawanan perpindahan; kerjanya \emph{negatif}.
\item \textbf{Benar.} Teorema usaha-energi: $\Delta K=W_{\text{net}}=F_{\text{net}}\cdot s$ (konstan untuk gaya konstan).
\item \textbf{Salah.} Saat $\alpha$ berubah, komponen vertikal $F$ berubah, mengubah gaya normal dan gaya gesek.
\item \textbf{Benar.} Perpindahan horizontal; gravitasi dan normal tegak lurus perpindahan, kerjanya nol.
\end{enumerate}
\end{pembox}

\begin{soal}{7}
Gelombang berdiri ujung terikat: $y=0{,}2\sin(5\pi x)\cos(2\pi t)$
($x,y$ dalam cm; $t$ dalam s). Tentukan kebenaran setiap pernyataan!

\par\smallskip
{\renewcommand{\arraystretch}{1.4}%
\begin{tabular}{@{}p{0.76\linewidth}cc@{}}
\hline\textbf{Pernyataan}&\textbf{Benar}&\textbf{Salah}\\\hline
(1)~Periode gelombang adalah 1 sekon.&$\bigcirc$&$\bigcirc$\\
(2)~Amplitudo gelombang bernilai konstan, yaitu 0{,}2 cm.&$\bigcirc$&$\bigcirc$\\
(3)~Pada jarak 0{,}4 cm dari ujung terikat terbentuk simpul.&$\bigcirc$&$\bigcirc$\\
(4)~Cepat rambat gelombang adalah 0{,}4 m/s.&$\bigcirc$&$\bigcirc$\\\hline
\end{tabular}}
\end{soal}
\begin{pembox}
\jawaban{(1)~\textbf{B}\quad(2)~\textbf{S}\quad(3)~\textbf{B}\quad(4)~\textbf{S}}
$k=5\pi$ rad/cm, $\omega=2\pi$ rad/s.
\begin{enumerate}[leftmargin=*,topsep=2pt,itemsep=2pt]
\item \textbf{Benar.} $T=2\pi/\omega=1$ s.
\item \textbf{Salah.} Amplitudo lokal $=0{,}2|\sin(5\pi x)|$, bergantung posisi (bukan konstan).
\item \textbf{Benar.} Simpul saat $\sin(5\pi x)=0\Rightarrow x=n/5$ cm. Untuk $x=0{,}4$ cm: $n=2$ ✓.
\item \textbf{Salah.} $\lambda=2\pi/(5\pi)=0{,}4$ cm; $v=f\lambda=1\times0{,}4$ cm/s $=0{,}004$ m/s, bukan 0{,}4 m/s.
\end{enumerate}
\end{pembox}

\begin{soal}{8}
\gambar[0.42\textwidth]{images/4/4-n8.png}
Balok $A$ (100 N) di atas balok $B$ (500 N), $A$ diikat tali. Koefisien gesek
$A$--$B$: statik 0{,}2, kinetik 0{,}1; $B$--lantai: statik 0{,}6, kinetik 0{,}5.
Gaya $F$ minimum untuk menggeser $B$ adalah \ldots\ N.
\begin{pilB}
  \item $380$ \item $400$ \item $420$ \item $440$ \item $460$
\end{pilB}
\end{soal}
\begin{pembox}
\jawaban{(a) 380 N}
Karena $A$ terikat dinding, saat $B$ bergerak gesek $A$--$B$ berarah ke kiri.
Normal lantai $=100+500=600$ N; gesek statik lantai maks: $0{,}6\times600=360$ N.
Gesek statik maks $A$--$B$: $0{,}2\times100=20$ N (berarah menahan $B$).
$F_{\min}=360+20=380$ N.
\end{pembox}

\begin{soal}{9}
\gambar[0.4\textwidth]{images/4/4-n9.png}
Dua vektor $A$ dan $B$ (20 N) di bidang $xy$ seperti gambar. Komponen
resultan pada sumbu $x$ nol. Besar vektor $A$ adalah \ldots\ N.
\begin{pilB}
  \item $6$ \item $8$ \item $10$ \item $12$ \item $14$
\end{pilB}
\end{soal}
\begin{pembox}
\jawaban{(c) 10 N}
\konsep $R_x=0\Rightarrow A_x+B_x=0$.

Dari gambar, vektor $A$ horizontal (ke kiri, $A_x=-A$) dan vektor $B$ (20 N)
membentuk sudut $60^\circ$ terhadap sumbu $x$ (ke kanan-atas, $B_x=20\cos60^\circ=10$ N).
Syarat $R_x=0$: $-A+10=0\Rightarrow A=10$ N.
\begin{catatbox}
Nilai sudut diambil dari gambar. Tanpa gambar, besar $A$ tidak dapat ditentukan.
\end{catatbox}
\end{pembox}

\begin{soal}{10}
$x(t)=6t^2-4t+2$ (m, s). Tentukan kebenaran setiap pernyataan!

\par\smallskip
{\renewcommand{\arraystretch}{1.4}%
\begin{tabular}{@{}p{0.76\linewidth}cc@{}}
\hline\textbf{Pernyataan}&\textbf{Benar}&\textbf{Salah}\\\hline
(1)~$v(t)=12t+4$ m/s.&$\bigcirc$&$\bigcirc$\\
(2)~Percepatan $a=12$ m/s$^2$ (konstan).&$\bigcirc$&$\bigcirc$\\
(3)~Kecepatan pada $t=3$ s: $v(3)=32$ m/s.&$\bigcirc$&$\bigcirc$\\
(4)~Kecepatan awal $v(0)=-4$ m/s.&$\bigcirc$&$\bigcirc$\\\hline
\end{tabular}}
\end{soal}
\begin{pembox}
\jawaban{(1)~\textbf{S}\quad(2)~\textbf{B}\quad(3)~\textbf{B}\quad(4)~\textbf{B}}
$v(t)=\dfrac{dx}{dt}=12t-4$ m/s;\quad$a=12$ m/s$^2$.
\begin{enumerate}[leftmargin=*,topsep=2pt,itemsep=2pt]
\item \textbf{Salah.} $v(t)=12t-4$, bukan $12t+4$.
\item \textbf{Benar.} $a=dv/dt=12$ m/s$^2$ (konstan).
\item \textbf{Benar.} $v(3)=36-4=32$ m/s.
\item \textbf{Benar.} $v(0)=0-4=-4$ m/s.
\end{enumerate}
\end{pembox}

% ===================================================================
%  NOMOR 11--20
% ===================================================================

\begin{soal}{11}
Partikel $Z$: energi total 17{,}0 GeV, massa diam $8{,}0$ GeV/$c^2$.
Energi kinetiknya \ldots\ GeV.
\begin{pilB}
  \item $17{,}0$ \item $25{,}0$ \item $136{,}0$ \item $8{,}0$ \item $9{,}0$
\end{pilB}
\end{soal}
\begin{pembox}
\jawaban{(e) 9,0 GeV}
\konsep $E=E_0+K\Rightarrow K=E-E_0$.

$K=17{,}0-8{,}0=9{,}0$ GeV.
\end{pembox}

\begin{soal}{12}
Dua bola identik dilempar horizontal ke dinding licin dengan kelajuan sama.
$P$ memantul balik sama kelajuan; $Q$ berhenti. $I_P/I_Q$ dan $F_P/F_Q$ = \ldots
\begin{pil}
  \item $I_P/I_Q=1$ dan $F_P/F_Q=1$
  \item $I_P/I_Q=1$ dan $F_P/F_Q=2$
  \item $I_P/I_Q=2$ dan $F_P/F_Q=1$
  \item $I_P/I_Q=2$ dan $F_P/F_Q=2$
  \item $I_P/I_Q=2$ dan $F_P/F_Q=\Delta t_Q/\Delta t_P$
\end{pil}
\end{soal}
\begin{pembox}
\jawaban{Tidak ada opsi yang sepenuhnya sesuai; $I_P/I_Q=2$ dan $F_P/F_Q=2\Delta t_Q/\Delta t_P$}
$I=|\Delta p|$. Bola $P$: $I_P=|-mv-mv|=2mv$. Bola $Q$: $I_Q=mv$.
$I_P/I_Q=2$. Gaya rata-rata $F=I/\Delta t$:
$\dfrac{F_P}{F_Q}=\dfrac{2mv/\Delta t_P}{mv/\Delta t_Q}=2\dfrac{\Delta t_Q}{\Delta t_P}$.
\begin{catatbox}
Opsi (e) memuat $I_P/I_Q=2$ (benar) tetapi $F_P/F_Q=\Delta t_Q/\Delta t_P$ (kurang faktor 2). Tidak ada opsi yang tepat.
\end{catatbox}
\end{pembox}

\begin{soal}{13}
\gambar[0.26\textwidth]{images/4/4-3.png}
Bola karet (0{,}20 kg) dijatuhkan dari 1{,}8 m di atas pegas. Gaya maks aman
36 N; $g=10$ m/s$^2$. Konstanta pegas minimum agar aman \ldots\ N/m.
\begin{pilB}
  \item $80$ \item $120$ \item $160$ \item $200$ \item $240$
\end{pilB}
\end{soal}
\begin{pembox}
\jawaban{(c) 160 N/m}
Kompresi maks $x$; bola turun $h+x$: $mg(h+x)=\tfrac12kx^2$.
Gaya maks $F=kx=36$ N $\Rightarrow x=36/k$.
Substitusi: $mg\!\left(h+\dfrac{36}{k}\right)=\dfrac{1}{2}k\!\left(\dfrac{36}{k}\right)^2=\dfrac{36^2}{2k}$.

Kalikan $k$: $mghk+36mg=\dfrac{36^2}{2}=648$.
\[
  k=\frac{648-36mg}{mgh}=\frac{648-36(2)}{2(1{,}8)}=\frac{648-72}{3{,}6}=\frac{576}{3{,}6}=160\text{ N/m}.
\]
\end{pembox}

\begin{soal}{14}
\gambar[0.32\textwidth]{images/4/4-4.png}
$P$ (0{,}30 kg, ketinggian 20 cm) menumbuk $Q$ (0{,}20 kg, diam) dan menempel.
Tentukan kebenaran setiap pernyataan!

\par\smallskip
{\renewcommand{\arraystretch}{1.4}%
\begin{tabular}{@{}p{0.76\linewidth}cc@{}}
\hline\textbf{Pernyataan}&\textbf{Benar}&\textbf{Salah}\\\hline
(1)~Kecepatan $P$ di titik terendah: $v=\sqrt{2gH}=2$ m/s.&$\bigcirc$&$\bigcirc$\\
(2)~Setelah menempel: $v'=\dfrac{m_P}{m_P+m_Q}v=0{,}6v=1{,}2$ m/s.&$\bigcirc$&$\bigcirc$\\
(3)~Tinggi maksimum setelah tumbukan: $h'=0{,}6H=12$ cm.&$\bigcirc$&$\bigcirc$\\
(4)~Tinggi maksimum setelah tumbukan: $h'=0{,}36H=7{,}2$ cm.&$\bigcirc$&$\bigcirc$\\\hline
\end{tabular}}
\end{soal}
\begin{pembox}
\jawaban{(1)~\textbf{B}\quad(2)~\textbf{B}\quad(3)~\textbf{S}\quad(4)~\textbf{B}}
\begin{enumerate}[leftmargin=*,topsep=2pt,itemsep=2pt]
\item \textbf{Benar.} $v=\sqrt{2\times10\times0{,}20}=2$ m/s.
\item \textbf{Benar.} $v'=\tfrac{0{,}30}{0{,}50}\times2=1{,}2$ m/s.
\item \textbf{Salah.} $h'=\dfrac{v'^2}{2g}=\dfrac{(1{,}2)^2}{20}=0{,}072$ m $=7{,}2$ cm.
      Pernyataan (3) keliru: $h'=(v'/v)\cdot H$ tidak benar; seharusnya $(v'/v)^2\cdot H$.
\item \textbf{Benar.} $h'=0{,}36\times20=7{,}2$ cm.
\end{enumerate}
\end{pembox}

\begin{soal}{15}
Planet $X$: jari-jari $1{,}5R_B$, massa jenis $2\rho_B$. Berat benda di Bumi $W$.
Berat benda di planet $X$ \ldots
\begin{pilB}
  \item $\tfrac{2}{3}W$ \item $W$ \item $\tfrac{3}{2}W$ \item $2W$ \item $3W$
\end{pilB}
\end{soal}
\begin{pembox}
\jawaban{(e) $3W$}
$g=\dfrac{GM}{R^2}$, $M=\rho\cdot\tfrac43\pi R^3$, sehingga $g\propto\rho R$.
$\dfrac{g_X}{g_B}=\dfrac{\rho_X R_X}{\rho_B R_B}=2\times1{,}5=3$.
Jadi $W_X=3W$.
\begin{catatbox}
Opsi (e) $3W$ merupakan pilihan yang diupdate; soal asli SM-SSU 2025 tidak memuat $3W$.
\end{catatbox}
\end{pembox}

\begin{soal}{16}
\gambar[0.5\textwidth]{images/4/4-5.png}
Batang homogen panjang 4 m berat 100 N ditumpu di kedua ujungnya. Anak
300 N berdiri 1 m dari kiri; beban 200 N digantung 0{,}5 m dari kanan.
Gaya tumpu ujung kanan adalah \ldots\ N.

\vspace{0.4em}\noindent Jawaban:\enspace\underline{\hspace{4cm}}
\end{soal}
\begin{pembox}
\jawaban{$R_B = 300$ N}
Momen terhadap ujung kiri ($A$); $R_A$ tidak masuk:
\[
  R_B\times4=300\times1+100\times2+200\times3{,}5=300+200+700=1200.
\]
$R_B=300$ N.
\end{pembox}

\begin{soal}{17}
\gambar[0.42\textwidth]{images/4/4-6.png}
Dua pegas identik $k$ tersusun seri; ditarik 20 cm; usaha 4 J.
Tentukan kebenaran setiap pernyataan!

\par\smallskip
{\renewcommand{\arraystretch}{1.4}%
\begin{tabular}{@{}p{0.76\linewidth}cc@{}}
\hline\textbf{Pernyataan}&\textbf{Benar}&\textbf{Salah}\\\hline
(1)~$k_{\text{eq}}=\dfrac{k}{2}$ untuk dua pegas seri identik.&$\bigcirc$&$\bigcirc$\\
(2)~Dari data: $W=4$ J, $x=0{,}20$ m, diperoleh $k=200$ N/m.&$\bigcirc$&$\bigcirc$\\
(3)~Nilai $k$ setiap pegas adalah 400 N/m.&$\bigcirc$&$\bigcirc$\\
(4)~Jika disusun paralel, $k_{\text{eq}}=2k=800$ N/m.&$\bigcirc$&$\bigcirc$\\\hline
\end{tabular}}
\end{soal}
\begin{pembox}
\jawaban{(1)~\textbf{B}\quad(2)~\textbf{S}\quad(3)~\textbf{B}\quad(4)~\textbf{B}}
$k_{\text{eq}}=k/2$.
$W=\tfrac12 k_{\text{eq}}x^2=\tfrac12\cdot\tfrac{k}{2}\cdot(0{,}20)^2=0{,}01k=4\Rightarrow k=400$ N/m.
\begin{enumerate}[leftmargin=*,topsep=2pt,itemsep=2pt]
\item \textbf{Benar.} $k_{\text{eq}}=k/2$.
\item \textbf{Salah.} Perhitungan yang benar: $0{,}01k=4\Rightarrow k=400$ N/m, bukan 200 N/m.
\item \textbf{Benar.} $k=400$ N/m.
\item \textbf{Benar.} Paralel: $k_{\text{eq}}=k+k=800$ N/m.
\end{enumerate}
\end{pembox}

\begin{soal}{18}
\gambar[0.42\textwidth]{images/4/4-n18.png}
Grafik $F$-$t$ segitiga: puncak 160 N, lama 0{,}010 s. Bola (0{,}05 kg)
bergerak mendekati dinding 12 m/s, berbalik arah. Kelajuan sesudah \ldots\ m/s.

\vspace{0.4em}\noindent Jawaban:\enspace\underline{\hspace{4cm}}
\end{soal}
\begin{pembox}
\jawaban{$4$ m/s}
Impuls $I=\tfrac12\times0{,}010\times160=0{,}8$ N\,s.
$\dfrac{I}{m}=\dfrac{0{,}8}{0{,}05}=16$ m/s (perubahan kecepatan).
Ambil arah menuju dinding positif; impuls dinding berlawanan:
$v_f=12-16=-4$ m/s. Kelajuan $=4$ m/s.
\end{pembox}

\begin{soal}{19}
$U=\tfrac12kx^2$; saat $x=10$ cm, $U=0{,}80$ J. Massa 0{,}50 kg berosilasi.
Tentukan kebenaran setiap pernyataan!

\par\smallskip
{\renewcommand{\arraystretch}{1.4}%
\begin{tabular}{@{}p{0.76\linewidth}cc@{}}
\hline\textbf{Pernyataan}&\textbf{Benar}&\textbf{Salah}\\\hline
(1)~$k=\dfrac{2U}{x^2}=160$ N/m.&$\bigcirc$&$\bigcirc$\\
(2)~Periode: $T=2\pi\sqrt{m/k}$.&$\bigcirc$&$\bigcirc$\\
(3)~Dengan $m=0{,}50$ kg, $k=160$ N/m: $T\approx0{,}50$ s.&$\bigcirc$&$\bigcirc$\\
(4)~Dengan $m=0{,}50$ kg, $k=160$ N/m: $T\approx0{,}35$ s.&$\bigcirc$&$\bigcirc$\\\hline
\end{tabular}}
\end{soal}
\begin{pembox}
\jawaban{(1)~\textbf{B}\quad(2)~\textbf{B}\quad(3)~\textbf{S}\quad(4)~\textbf{B}}
$k=\dfrac{2(0{,}80)}{(0{,}10)^2}=160$ N/m.
$T=2\pi\sqrt{\dfrac{0{,}50}{160}}=2\pi\times0{,}05590\approx0{,}351$ s.
\begin{enumerate}[leftmargin=*,topsep=2pt,itemsep=2pt]
\item \textbf{Benar.} $k=160$ N/m.
\item \textbf{Benar.} Formula periode benar.
\item \textbf{Salah.} $T\approx0{,}35$ s, bukan 0{,}50 s.
\item \textbf{Benar.} $T\approx0{,}35$ s.
\end{enumerate}
\end{pembox}

\begin{soal}{20}
$y_1=4\sin(10t)$ cm, $y_2=3\cos(10t)$ cm. Tentukan kebenaran setiap pernyataan!

\par\smallskip
{\renewcommand{\arraystretch}{1.4}%
\begin{tabular}{@{}p{0.76\linewidth}cc@{}}
\hline\textbf{Pernyataan}&\textbf{Benar}&\textbf{Salah}\\\hline
(1)~Amplitudo resultan $y_1+y_2$ adalah $4+3=7$ cm.&$\bigcirc$&$\bigcirc$\\
(2)~$y_2=3\sin(10t+90^\circ)$; beda fase $y_1$ dan $y_2$ adalah $90^\circ$.&$\bigcirc$&$\bigcirc$\\
(3)~Amplitudo resultan $R=\sqrt{4^2+3^2}=5$ cm.&$\bigcirc$&$\bigcirc$\\
(4)~Kedua getaran berfrekuensi sudut 10 rad/s.&$\bigcirc$&$\bigcirc$\\\hline
\end{tabular}}
\end{soal}
\begin{pembox}
\jawaban{(1)~\textbf{S}\quad(2)~\textbf{B}\quad(3)~\textbf{B}\quad(4)~\textbf{B}}
\begin{enumerate}[leftmargin=*,topsep=2pt,itemsep=2pt]
\item \textbf{Salah.} Amplitudo tidak bisa langsung dijumlahkan bila berbeda fase; $R=\sqrt{A_1^2+A_2^2}=5$ cm.
\item \textbf{Benar.} $\cos\theta=\sin(\theta+90^\circ)$, beda fase $90^\circ$.
\item \textbf{Benar.} $R=5$ cm.
\item \textbf{Benar.} Keduanya $\omega=10$ rad/s.
\end{enumerate}
\end{pembox}

% ===================================================================
%  NOMOR 21--30
% ===================================================================

\begin{soal}{21}
\gambar[0.4\textwidth]{images/4/4-7.png}
GHS: amplitudo 6 cm, $f=0{,}5$ Hz. Pada $t=0$ partikel di titik setimbang
bergerak ke kiri. Posisi pada $t=0{,}5$ s adalah $x=\ldots$ cm.
\begin{pilB}
  \item $-6$ \item $-3$ \item $0$ \item $3$ \item $6$
\end{pilB}
\end{soal}
\begin{pembox}
\jawaban{(a) $x=-6$ cm}
$T=1/f=2$ s; $t=0{,}5$ s $=T/4$.

Pada $t=0$: $x=0$, bergerak ke kiri. Setelah $T/4$, partikel mencapai
simpangan maksimum ke arah geraknya, yaitu $x=-A=-6$ cm.
\end{pembox}

\begin{soal}{22}
Pipa organa terbuka $L=85$ cm, $v_s=340$ m/s. Frekuensi nada atas kedua \ldots\ Hz.

\vspace{0.4em}\noindent Jawaban:\enspace\underline{\hspace{4cm}}
\end{soal}
\begin{pembox}
\jawaban{$600$ Hz}
Pipa terbuka: $f_n=n\dfrac{v}{2L}$. Frekuensi dasar ($n=1$):
$f_1=\dfrac{340}{2\times0{,}85}=200$ Hz.
Nada atas kedua adalah harmonik ketiga ($n=3$):
$f_3=3\times200=600$ Hz.
\end{pembox}

\begin{soal}{23}
$y=0{,}04\sin(\pi x-20\pi t)$ (m, s). Massa tali 0{,}80 g; panjang 2{,}0 m.
Tegangan tali \ldots\ N.

\vspace{0.4em}\noindent Jawaban:\enspace\underline{\hspace{4cm}}
\end{soal}
\begin{pembox}
\jawaban{$0{,}16$ N}
$k=\pi$ rad/m, $\omega=20\pi$ rad/s.
$v=\omega/k=20$ m/s.
$\mu=m/L=0{,}80\times10^{-3}/2{,}0=4{,}0\times10^{-4}$ kg/m.
$T=\mu v^2=4{,}0\times10^{-4}\times400=0{,}16$ N.
\end{pembox}

\begin{soal}{24}
Gelombang transversal ke kiri: $A=5$ cm, $\lambda=2$ m, $f=4$ Hz.
Pada $x=0$, $t=0$: $y=0$, medium bergerak ke atas.
Tentukan kebenaran setiap pernyataan!

\par\smallskip
{\renewcommand{\arraystretch}{1.4}%
\begin{tabular}{@{}p{0.76\linewidth}cc@{}}
\hline\textbf{Pernyataan}&\textbf{Benar}&\textbf{Salah}\\\hline
(1)~$k=\pi$ rad/m dan $\omega=8\pi$ rad/s.&$\bigcirc$&$\bigcirc$\\
(2)~Gelombang ke kiri: tanda positif pada $\omega t$.&$\bigcirc$&$\bigcirc$\\
(3)~Pada $x=0$, $t=0$: $y=0{,}05\sin(0)=0$ sesuai syarat.&$\bigcirc$&$\bigcirc$\\
(4)~Pada $x=0$, $t=0$, medium bergerak ke bawah.&$\bigcirc$&$\bigcirc$\\\hline
\end{tabular}}
\end{soal}
\begin{pembox}
\jawaban{(1)~\textbf{B}\quad(2)~\textbf{B}\quad(3)~\textbf{B}\quad(4)~\textbf{S}}
Persamaan: $y=0{,}05\sin(\pi x+8\pi t)$.
\begin{enumerate}[leftmargin=*,topsep=2pt,itemsep=2pt]
\item \textbf{Benar.} $k=2\pi/2=\pi$; $\omega=2\pi(4)=8\pi$.
\item \textbf{Benar.} Gelombang ke $-x$: $(kx+\omega t)$.
\item \textbf{Benar.} $\sin(0)=0$ ✓.
\item \textbf{Salah.} $\dfrac{\partial y}{\partial t}\big|_{(0,0)}=0{,}05\times8\pi\times\cos(0)>0$ (ke \emph{atas}).
\end{enumerate}
\end{pembox}

\begin{soal}{25}
$y=10\sin(4\pi x)\cos(50\pi t)$ ($x$ dalam m, $y$ dalam cm).
Jarak antara dua simpul berurutan \ldots\ m.

\vspace{0.4em}\noindent Jawaban:\enspace\underline{\hspace{4cm}}
\end{soal}
\begin{pembox}
\jawaban{$0{,}25$ m}
Simpul saat $\sin(4\pi x)=0\Rightarrow4\pi x=n\pi\Rightarrow x=n/4$ m.
Dua simpul berurutan ($n$ dan $n+1$): $\Delta x=1/4=0{,}25$ m.
\end{pembox}

\begin{soal}{26}
\gambar[0.5\textwidth]{images/4/4-8.png}
Jarak puncak 1 ke puncak 3 adalah 1{,}2 m; 8 gelombang dalam 4 s.
Tentukan kebenaran setiap pernyataan!

\par\smallskip
{\renewcommand{\arraystretch}{1.4}%
\begin{tabular}{@{}p{0.76\linewidth}cc@{}}
\hline\textbf{Pernyataan}&\textbf{Benar}&\textbf{Salah}\\\hline
(1)~Jarak puncak 1 ke puncak 3 = $2\lambda$.&$\bigcirc$&$\bigcirc$\\
(2)~$2\lambda=1{,}2$ m $\Rightarrow\lambda=0{,}6$ m.&$\bigcirc$&$\bigcirc$\\
(3)~8 gelombang dalam 4 s: $f=2$ Hz.&$\bigcirc$&$\bigcirc$\\
(4)~$v=f\lambda=2\times0{,}6=0{,}6$ m/s.&$\bigcirc$&$\bigcirc$\\\hline
\end{tabular}}
\end{soal}
\begin{pembox}
\jawaban{(1)~\textbf{B}\quad(2)~\textbf{B}\quad(3)~\textbf{B}\quad(4)~\textbf{S}}
\begin{enumerate}[leftmargin=*,topsep=2pt,itemsep=2pt]
\item \textbf{Benar.} Puncak ke-1 ke puncak ke-3 melewati 2 periode ruang.
\item \textbf{Benar.} $\lambda=0{,}6$ m.
\item \textbf{Benar.} $f=8/4=2$ Hz.
\item \textbf{Salah.} $v=2\times0{,}6=1{,}2$ m/s, bukan 0{,}6 m/s.
\end{enumerate}
\end{pembox}

\begin{soal}{27}
$y_1=3\cos(kx-\omega t)$, $y_2=3\cos(kx+\omega t)$.
Tentukan kebenaran setiap pernyataan!

\par\smallskip
{\renewcommand{\arraystretch}{1.4}%
\begin{tabular}{@{}p{0.76\linewidth}cc@{}}
\hline\textbf{Pernyataan}&\textbf{Benar}&\textbf{Salah}\\\hline
(1)~Resultan merupakan gelombang stasioner.&$\bigcirc$&$\bigcirc$\\
(2)~Amplitudo maksimum resultan adalah 6.&$\bigcirc$&$\bigcirc$\\
(3)~Simpul saat $kx=\dfrac{(2n+1)\pi}{2}$, $n\in\mathbb{Z}$.&$\bigcirc$&$\bigcirc$\\
(4)~Kedua gelombang memiliki frekuensi yang sama.&$\bigcirc$&$\bigcirc$\\\hline
\end{tabular}}
\end{soal}
\begin{pembox}
\jawaban{(1)~\textbf{B}\quad(2)~\textbf{B}\quad(3)~\textbf{B}\quad(4)~\textbf{B}}
$y_1+y_2=2\cdot3\cos(kx)\cos(\omega t)=6\cos(kx)\cos(\omega t)$ (gelombang stasioner).
\begin{enumerate}[leftmargin=*,topsep=2pt,itemsep=2pt]
\item \textbf{Benar.} Bentuk $f(x)\cdot g(t)$ adalah gelombang stasioner.
\item \textbf{Benar.} Amplitudo maks $=6$ (saat $|\cos(kx)|=1$).
\item \textbf{Benar.} Simpul saat $\cos(kx)=0\Rightarrow kx=(2n+1)\pi/2$.
\item \textbf{Benar.} $\omega$ sama $\Rightarrow f$ sama.
\end{enumerate}
\end{pembox}

\begin{soal}{28}
$y_1=4\cos(kx-\omega t)$, $y_2=6\cos(kx+\omega t)$.
Amplitudo maks dan min resultan berturut-turut \ldots
\begin{pilB}
  \item $2$ dan $10$ \item $4$ dan $6$ \item $6$ dan $4$
  \item $10$ dan $2$ \item $12$ dan $0$
\end{pilB}
\end{soal}
\begin{pembox}
\jawaban{(d) 10 dan 2}
Dua gelombang sama frekuensi berlawanan arah; amplitudo resultan bervariasi terhadap posisi.
$A_{\max}=A_1+A_2=4+6=10$;\quad$A_{\min}=|A_2-A_1|=|6-4|=2$.
\end{pembox}

\begin{soal}{29}
Jendela 0{,}40 m$^2$, bunyi 80 dB ($I_0=10^{-12}$ W/m$^2$).
Daya akustik masuk $=4\times\ldots$ W.

\vspace{0.4em}\noindent Jawaban:\enspace\underline{\hspace{4cm}}
\end{soal}
\begin{pembox}
\jawaban{$10^{-5}$}
$80=10\log(I/10^{-12})\Rightarrow I=10^{-4}$ W/m$^2$.
$P=IA=10^{-4}\times0{,}40=4{,}0\times10^{-5}$ W $=4\times10^{-5}$ W.
Bagian yang kosong: $10^{-5}$.
\end{pembox}

\begin{soal}{30}
Speaker titik, daya tetap. Taraf intensitas di jarak 2 m: 70 dB.
Tentukan kebenaran setiap pernyataan!

\par\smallskip
{\renewcommand{\arraystretch}{1.4}%
\begin{tabular}{@{}p{0.76\linewidth}cc@{}}
\hline\textbf{Pernyataan}&\textbf{Benar}&\textbf{Salah}\\\hline
(1)~Untuk sumber titik: $I\propto1/r^2$.&$\bigcirc$&$\bigcirc$\\
(2)~Jarak bertambah 10$\times$: intensitas berkurang 10$\times$ ($\Delta\beta=-10$ dB).&$\bigcirc$&$\bigcirc$\\
(3)~$\Delta\beta=10\log(1/100)=-20$ dB.&$\bigcirc$&$\bigcirc$\\
(4)~Taraf intensitas pada jarak 20 m: $\beta_2=50$ dB.&$\bigcirc$&$\bigcirc$\\\hline
\end{tabular}}
\end{soal}
\begin{pembox}
\jawaban{(1)~\textbf{B}\quad(2)~\textbf{S}\quad(3)~\textbf{B}\quad(4)~\textbf{B}}
\begin{enumerate}[leftmargin=*,topsep=2pt,itemsep=2pt]
\item \textbf{Benar.} Sumber titik: $I\propto1/r^2$.
\item \textbf{Salah.} Jarak naik 10$\times\Rightarrow I$ turun $100\times$ (bukan 10$\times$), sehingga $\Delta\beta=-20$ dB.
\item \textbf{Benar.} $\Delta\beta=10\log(1/100)=-20$ dB.
\item \textbf{Benar.} $\beta_2=70-20=50$ dB.
\end{enumerate}
\end{pembox}

\end{document}
